\label{Conclusions}

This project set out to determine whether combining residual learning with the UNet backbone
conventionally used as the denoising network of a diffusion model improves the quality of the images
it generates. A Denoising Diffusion Probabilistic Model was implemented from first principles and
the diffusion framework was held fixed while three denoising networks were substituted into it: the
proposed residual UNet, an otherwise identical network with the identity shortcut removed, and the
reference architecture of Ho et al.\ narrowed to a matching parameter budget. Each was trained on
CIFAR-10 at $32 \times 32$ and on CelebA at $64 \times 64$ under identical conditions.

\section*{Findings}

The answer to the question as posed is that residual learning helps, but conditionally, and not in
the way the project anticipated.

The shortcut makes optimisation easier at the outset. After one epoch the proposed network leads its
shortcut-free twin in all six paired runs, by $5.4\%$ on average on CIFAR-10 and $11.1\%$ on
CelebA, and since the two differ in nothing else the shortcut is the only possible cause. That
advantage is gone by the fifth epoch and does not reappear: at convergence the restoration measures
cannot separate the two networks, and the two that lean at all lean in opposite directions, the
training loss favouring the plain network and structural similarity the proposed one, both by less
than half of one per cent.

They are not, however, indistinguishable as generators, and repeating each configuration under three
seeds is what makes the distinction legible. On CelebA the proposed network leads in every seed, by
$59$, $22$ and $109$ Fréchet Inception Distance, for a mean advantage of $63.31$ against a
seed-to-seed standard deviation of about $23$ for either network. On CIFAR-10 the same comparison
yields $-3.46 \pm 13.78$ with the sign changing between seeds, which is to say no effect at all;
the apparent reversal seen in a single run was noise. The difference between the two datasets tracks
depth, the CelebA configuration having three resolution stages against two for CIFAR-10, and
residual connections exist precisely to make deeper networks optimisable. Where the network is
shallow enough not to need one, adding it changes nothing; where it is deep enough to struggle
without one, it is worth sixty points of FID.

The reference architecture remains the strongest network overall. It leads on every restoration
measure on both datasets and on CIFAR-10 FID by a wide margin, while using fewer parameters than
either competitor, at roughly $2.3$ times the training cost per epoch. Only on CelebA FID is it
matched, and there the proposed network achieves the same quality for less than half the compute.

Two results of a different kind emerged from the work and are worth recording. The first is
methodological: a network can denoise as well as another and generate images that are $78\%$ worse
by FID. The training objective averages the noise-prediction error uniformly over timesteps, whereas
generation composes a thousand such predictions in sequence, so errors that are small on average but
systematic at particular timesteps compound through that composition. Evaluating a diffusion model
by reconstruction alone is not sufficient. The second is practical: the algebraic form of the
posterior mean used in the reverse process is numerically stable only for a network whose noise
prediction is already accurate. Without clipping the implied estimate of $x_0$ at every step, the
reverse process diverged from a standard deviation of $1$ to approximately $16$, in single precision
as well as mixed, and produced saturated noise from every model.

\section*{Limitations}

Each configuration was trained three times. Three runs are enough to establish that a sign is
consistent and to estimate a standard deviation, but not to support a formal claim of significance:
the paired test on the CelebA result gives $t(2) = -2.50$, which does not reach the conventional
threshold with two degrees of freedom. The finding rests on the direction being the same in every
run and the magnitude exceeding the observed spread, not on a $p$-value. The depth explanation
offered for the difference between datasets is consistent with the evidence but not established by
it, since depth and resolution vary together here; separating them would require the two-stage and
three-stage configurations to be trained on the same images.

The seed study also showed that the Fréchet Inception Distance is far noisier at this scale than the
restoration measures, with standard deviations between $7.8$ and $32.9$ against thousandths for
SSIM. Any comparison of diffusion models drawn from a single training run of this length, including
several reported in the earlier sections of this work before the repetition was carried out, should
be treated with corresponding caution.

Thirty epochs amounts to roughly eleven thousand optimisation steps, some two orders of magnitude
short of the schedule used for the published CIFAR-10 configuration, and the networks are about
fourteen times smaller. The absolute FID values are correspondingly poor and should not be compared
with the literature; only the relative ordering between the three networks, which were trained
identically, carries meaning. The Fréchet Inception Distance was computed from two thousand samples,
which is the smallest number at which the estimate is not degenerate given the dimensionality of the
Inception feature space, and remains biased upward at that size.

One standard practice was omitted. Diffusion implementations customarily sample not from the trained weights themselves but from an exponential moving average of them, which suppresses the fluctuation of the final iterates and typically improves the Fr\'echet Inception Distance appreciably. No such average was maintained here, so the absolute sample quality reported is lower than the same training run would otherwise have produced. The omission applies identically to all six runs and so does not affect the comparison between them.

The memorisation check operates in pixel space and would not detect a reproduction that had been
shifted or recoloured. Negative Log-Likelihood was excluded on the grounds that this implementation
fixes the reverse-process variances rather than learning them. The comparison against Variational
Autoencoders and Generative Adversarial Networks set out in Chapter~\ref{ch2} was not carried out.

\section*{Further work}

The immediate priority is more repetition. Three seeds established which of the single-run findings
survive, but the surviving one is supported by consistency of sign rather than by a significance
test; ten seeds would permit a confidence interval, and at the cost of the deterministic sampler
described below would remain affordable.

The depth hypothesis is directly testable. Training two-stage, three-stage and four-stage variants
of the proposed network on a single dataset, each with and without the shortcut, would isolate depth
from resolution and determine whether the benefit of residual learning grows with the former as the
argument predicts.

A separate line of work concerns the absolute quality of the samples rather than the comparison between architectures. Three changes would address it, in descending order of expected effect and ascending order of cost. Maintaining an exponential moving average of the weights and sampling from it is the cheapest, requiring a few lines and no change to the experimental design. Extending training by an order of magnitude is the most certain, since Section~\ref{sec:indistinct} attributes the indistinctness principally to the budget. Replacing the ancestral sampler with the deterministic variant of Song et al.\ would allow comparable quality from roughly fifty reverse steps instead of a thousand, which would also reduce the cost of computing the Fr\'echet Inception Distance by a factor of twenty and so make the repeated runs recommended above considerably cheaper. None of the three alters the architectures under test, so each could be applied to all six configurations without disturbing the comparison.

Beyond that, learning the reverse-process variances rather than fixing them would make the
variational bound informative and bring Negative Log-Likelihood within reach, and the Variational
Autoencoder and Generative Adversarial Network baselines would place the diffusion results in the
wider context that Chapter~\ref{ch2} describes but does not measure.
